\documentclass[../../main.tex]{subfiles}
\begin{document}

\begin{flushright}
\footnotesize\textit{【自動文字起こし・要確認】}
\end{flushright}

{\bf [解]} $\alpha = \dfrac{-1}{2+4i}$ とすると，$\alpha$ は $\alpha = (3+4i)\alpha + 1$ をみたすので，
\begin{align}
z_{n+1} - \alpha &= (3+4i)(z_n - \alpha)
\end{align}
$z_1 = 1$ とあわせて，くり返し用いると
\begin{align}
z_n &= (3+4i)^{n-1}(1 - \alpha) + \alpha \\
&= \dfrac{1}{2+4i}(3+4i)^{n-1} - \dfrac{1}{2+4i} \\
&= \dfrac{1-2i}{10}\{(3+4i)^{n-1} - 1\} \quad \cdots \textcircled{1}
\end{align}
である．

\bigskip

\noindent
\paragraph{(1)}

与式の各辺 $0$ 以上だから，2乗して
\begin{align}
\dfrac{9}{16} \cdot 5^{2n-2} < |z_n|^2 < \dfrac{1}{16} \cdot 5^{2n} \quad \cdots \textcircled{2}
\end{align}
又，\textcircled{1}から
\begin{align}
|z_n| &= \left|\dfrac{1-2i}{10}\right| |(3+4i)^{n-1} - 1|
\end{align}
で，$\left|\dfrac{1-2i}{10}\right| = \dfrac{\sqrt{5}}{10}$ から，三角不等式より
\begin{align}
\begin{cases}
|z_n| \le \dfrac{\sqrt{5}}{10} \{ |3+4i|^{n-1} + |-1| \} = \dfrac{\sqrt{5}}{10} (5^{n-1} + 1) \\
|z_n| \ge \dfrac{\sqrt{5}}{10} \{ |3+4i|^{n-1} - |-1| \} = \dfrac{\sqrt{5}}{10} (5^{n-1} - 1)
\end{cases}
\end{align}
だから，$A = 5^n$ として，
\begin{align}
\dfrac{1}{20}(A-1)^2 \le |z_n|^2 \le \dfrac{1}{20}(A+1)^2 \quad \cdots \textcircled{3}
\end{align}
である．まず，\textcircled{2}の右側について，$A>0$ より
\begin{align}
\dfrac{1}{20}(A+1)^2 < \dfrac{1}{16} A^2 &\iff A^2 > 4(2A+1) \\
&\iff A > 8 + \dfrac{4}{A} \quad \cdots \textcircled{4}
\end{align}
であり，$n \ge 2$ の時，$A = 5^n \ge 25 > 8 + \dfrac{4}{25} \ge 8 + \dfrac{4}{A}$ から \textcircled{4} は成立．

\bigskip

\noindent
次に \textcircled{2}の左側について，
\begin{align}
\dfrac{1}{20}(A-1)^2 > \dfrac{9}{25 \cdot 16} A^2 &\iff 11A^2 - 40A + 20 > 0 \\
&\iff \left( A - \dfrac{20 + 6\sqrt{5}}{11} \right) \left( A + \dfrac{20 - 6\sqrt{5}}{11} \right) > 0 \quad \cdots \textcircled{5}
\end{align}
であり，$A \ge 25 > \dfrac{20 + 6\sqrt{5}}{11} > \dfrac{20 - 6\sqrt{5}}{11}$ から \textcircled{5} は常に成立．

\bigskip

\noindent
以上 \textcircled{4} 及び \textcircled{5} から，$n \ge 2$ の時，
\begin{align}
\dfrac{9}{16} \cdot 5^{2n-2} < \dfrac{1}{20}(A-1)^2 \le |z_n|^2 \le \dfrac{1}{20}(A+1)^2 < \dfrac{1}{16} \cdot 5^{2n}
\end{align}
で，両辺平方根をとって
\begin{align}
\dfrac{3}{4} \cdot 5^{n-1} < |z_n| < \dfrac{1}{4} \cdot 5^n
\end{align}
又，$n=1$ の時，$|z_1|=1$ から問題意は成立するから，以上より不等式は成立する． \hfill $\qed$

\bigskip

\noindent
\paragraph{(2)}

(1)から数列 $\{|z_n|\}$ は単調増加である．($\because \dfrac{1}{4} \cdot 5^{n-1} < \dfrac{3}{4} \cdot 5^n$) \\
したがって，$r = \dfrac{5^n}{4}$ の時 $f(r) = n$ であり，$\dfrac{5^{n-1}}{4} < r \le \dfrac{5^n}{4}$ に対して，$n-1 \le f(r) \le n$ が成立する．又この時
\begin{align}
(n-1)\log 5 - \log 4 < \log r \le n\log 5 - \log 4
\end{align}
が成立する．($\because \log r$ は単調増加) 以上の2式から，$n$ が十分大きい時
\begin{align}
\dfrac{n-1}{n\log 5 - \log 4} < \dfrac{f(r)}{\log r} < \dfrac{n}{(n-1)\log 5 - \log 4}
\end{align}
が成立する．$r \to \infty$ の時 $n \to \infty$ で，この両辺は共に $\dfrac{1}{\log 5}$ に収束するから，はさみうちより
\begin{align}
\dfrac{f(r)}{\log r} \longrightarrow \dfrac{1}{\log 5} \hfill \qed
\end{align}

\newpage

\end{document}
